\section{Introduction}\label{sec:introduction}

For a regular local domain $R$ with fraction field $F$, consider the
rational restriction map
\[
 K_j(R)_{\Q}\longrightarrow K_j(F)_{\Q},
 \qquad K_j(R)_{\Q}=K_j(R)\otimes_{\Z}\Q.
\]
Its injectivity is the first assertion of the augmented rational
Gersten conjecture. We construct classes in its kernel for ramified
regular local rings of mixed characteristic. The classes are represented
in actual algebraic $K$-groups: they have infinite order, and their
restrictions to the fraction field vanish before tensoring with $\Q$.

The two basic rings are
\begin{equation}\label{eq:intro-basic-rings}
\begin{aligned}
 A_2&=\left(\Z_{(5)}[T,x,y]/
       (\Phi_5(1+y)+x^4+Txy^3)\right)_{(5,x,y)},\\
 A_3&=\left(\Z_{(3)}[x,y,z]/(3+xy+z^2)\right)_{(x,y,z)},
\end{aligned}
\end{equation}
where $\Phi_5(t)=1+t+t^2+t^3+t^4$.
The ring $A_2$ has dimension two, maximal ideal $(x,y)$, and residue
field $\F_5(T)$; the ring $A_3$ has dimension three, maximal ideal
$(x,y,z)$, and residue field $\F_3$.
Write $\widehat A_i$ for completion at the maximal ideal.
Theorem~\ref{d2:thm:completed} and
Theorem~\ref{thm:explicit-three-dimensional} construct classes
$c_i\in K_3(A_i)$ whose images in $K_3(\widehat A_i)$ also have
infinite order. In each of the four rings the corresponding class
vanishes after inverting $x$, and hence in the fraction field,
integrally.

In dimension three the class is the product of three explicit units:
\begin{equation}\label{eq:intro-explicit-three}
 c_3=2\left\{\frac{1+z}{2},\,1+x,\,
                   1+\frac{(1+x)y}{4}\right\}.
\end{equation}
The braces denote the product of unit classes in Quillen $K_1$.
Dennis--Stein identities prove its integral vanishing after inverting
$x$. Its infinite order is detected on a smooth formal quadric by a
convergent three-adic endpoint series. The actual vanishing of the
special-fiber group $K_4$ makes the relative-to-absolute map injective,
so this argument detects an actual $K_3$-class.
The two-dimensional construction starts with a coefficient in the
$K_3$ of a cyclotomic integer ring and pushes it forward from the
regular divisor $x=0$. A second construction identifies the explicit
unit product
\[
 \left\{2+y,\ 1-x,\
 1+(1-x)(x^3+Ty^3)y\bigl((1+y)^5+1\bigr)\right\}
 \in K_3(A_2)
\]
as an infinite-order generic-kernel class in both $A_2$ and
$\widehat A_2$; see Proposition~\ref{d2ext:dimtwoExplicitTriple}.

The same methods give families with larger rank, higher $K$-degree,
and one fixed ring supporting infinitely many independent classes.
Table~\ref{tab:intro-results} locates these extensions. Every rank
assertion concerns the rational generic kernel; the witnesses are
actual integral classes vanishing on the indicated principal open.
The parameters introduced in residue fields preserve the local
Krull dimension.

\begin{table}[htbp]
\centering
\small
\renewcommand{\arraystretch}{1.18}
\begin{tabular}{@{}p{0.10\textwidth}p{0.64\textwidth}p{0.20\textwidth}@{}}
\hline
Dim. & Result & Reference\\
\hline
$2$ & At every residue prime, tame families have unbounded kernel rank
in degree three and nonzero kernel in every odd degree at least three.
& \ref{d2ext:twodimtamefamily}, \ref{d2ext:twodimtameallodd}\\
$3$ & Eisenstein families have explicit rank formulas before completion
and arbitrarily large lower bounds after completion.
& \ref{d3:thm:eisenstein}, \ref{d3:thm:completedeisensteinrank}\\
$3$ & Explicit cyclotomic triples at every prime retain their independence
after henselization and completion.
& \ref{d3:thm:explicitwildtriples}\\
$3$ & Products of $j$ specified units give infinite-order kernel classes
for every $j\geq3$.
& \ref{d3:thm:everydegreelaurent}\\
$2,3$ & At each prime, a fixed ring has infinite kernel rank in every odd
degree at least three; in dimension three this also holds for its completion.
& \ref{d2ext:twodimstrictcoefficientinfinite}, \ref{d3:thm:strictallprimeallodd}\\
$3$ & At every prime, one fixed ring and one fixed complete ring, with
residue field $\overline{\F}_p(T)$, have infinite kernel rank in every $j\geq3$.
& \ref{d3:thm:singlefixedalldegree}, \ref{d3:thm:completedfixedalldegree}\\
$2$ & One complete ring with residue field $\overline{\F}_5(T,S)$ has
infinite kernel rank in every $j\geq3$.
& \ref{d2ext:dimtwoCompleteAllDegreesInfiniteRank}\\
$2,3$ & Explicit degree-three classes yield Milnor kernels; the
three-dimensional class also gives Beilinson-motivic and higher-Chow kernels.
& \ref{d2ext:twodimexplicitMilnor}, \ref{d3:cor:explicitmilnormotivic}\\
\hline
\end{tabular}
\caption{Families and consequences of the basic constructions.
The cited statements give the coefficient rings and all rank bounds.}
\label{tab:intro-results}
\end{table}

The geometric and completed proofs detect the supported classes by
different cohomology theories. On an arithmetic quadric, a small
resolution shows that the ruling has infinite order in the local
class group. This calculation forces each geometric boundary component
removed by an allowed denominator to have zero divisor class.
Consequently the ruling's Betti class survives every finite
localization. Multiplying by a number-field coefficient with nonzero
regulator gives a nonzero analytic Deligne class on every such open.
Filtered continuity then proves nonvanishing on the local ring.
For the quartic surface, the corresponding divisor calculation uses
the exceptional quartic curve and its Jacobian.

Completion requires a separate detector. The Beilinson fiber square
of Antieau--Mathew--Morrow--Nikolaus identifies the relevant relative
$p$-completed $K$-theory with continuous cyclic homology
\cite{AMMN}. A finite mixed Hochschild--Kostant--Rosenberg comparison
retains the ordinary continuous de Rham quotients. On a quadric, an
explicit residue operator detects the ruling class. On the quartic
family, the elliptic quotients of its Jacobian and a bounded
slope-one functional detect the weighted root class after passage to
the generic Cohen completion. In each case the proof separately
controls the relative-to-absolute map and localization of the
integral test ring. This is the step that returns a cohomological
nonvanishing statement to the actual algebraic $K$-group.

For the rank constructions, cyclotomic classes supply finite families
of independent coefficient values. Their Chern classes factor through
spectral completion, so independence passes to the continuous cyclic
point coordinates without requiring surjectivity from actual
$K$-theory. A continuous Laurent constant-term operator preserves the
ruling detector after adjoining units. Products with these units raise
the $K$-degree, while unramified coefficient extensions increase the
rank. Finite descent handles each possible relation before completion;
a separate comparison for completed essentially smooth models handles
the fixed complete rings.

We use $K(R)$ for the algebraic $K$-theory spectrum and
$K(R;\Q_p)=K(R)^{\wedge}_p[1/p]$ for its rationalized spectral
$p$-completion. This notation is distinct from tensoring each actual
$K$-group with $\Q_p$. Throughout the continuous de Rham arguments,
cohomology means the quotient by the image of the differential, with
no closure of that image.
The final mathematical section gives an explicit obstruction to the
strict simplicial-functor construction in
\cite[Lemma~13(III)]{MochizukiParameters}.
The scope of the accompanying computational and formal checks is
recorded separately at the end of the article.
